% TODO make hw stylesheet
\documentclass[10pt]{scrartcl}
\usepackage[a4paper, total={6in, 8in}]{geometry}
\usepackage[usenames,dvipsnames,svgnames]{xcolor}
\usepackage[shortlabels]{enumitem}
\usepackage[framemethod=TikZ]{mdframed}
\usepackage{amsmath,amssymb,amsthm}
\usepackage[colorlinks]{hyperref}
\usepackage{multicol}
\usepackage{tabularx}

\hypersetup{
	colorlinks=true,
	linkcolor=blue,
	filecolor=magenta,      
	urlcolor=magenta,
	pdftitle={MAT 322 HW}
}

\usepackage{microtype}
\usepackage{mathtools}
\usepackage[headsepline]{scrlayer-scrpage}
\usepackage{thmtools}
\usepackage[textsize=footnotesize]{todonotes}

\mdfdefinestyle{mdbluebox}{roundcorner=10pt,innerbottommargin=9pt,
	linecolor=blue,backgroundcolor=TealBlue!5,}
\declaretheoremstyle[headfont=\sffamily\bfseries\color{MidnightBlue},
mdframed={style=mdbluebox},]{thmbluebox}

\mdfdefinestyle{mdbloobox}{frametitlefont=\bfseries,innerbottommargin=8pt,
	nobreak=true,backgroundcolor=TealBlue!5,linecolor=blue,}
\declaretheoremstyle[headfont=\bfseries\color{MidnightBlue},
mdframed={style=mdbloobox},headpunct={\\[3pt]},postheadspace=0pt,]{thmbloobox}

\mdfdefinestyle{mdredbox}{frametitlefont=\bfseries,innerbottommargin=8pt,
	nobreak=true,backgroundcolor=Salmon!5,linecolor=RawSienna,}
\declaretheoremstyle[headfont=\bfseries\color{RawSienna},
mdframed={style=mdredbox},headpunct={\\[3pt]},postheadspace=0pt,]{thmredbox}

\mdfdefinestyle{mdgreenbox}{linecolor=ForestGreen,backgroundcolor=ForestGreen!5,
	linewidth=2pt,rightline=false,leftline=true,topline=false,bottomline=false,}
\declaretheoremstyle[headfont=\bfseries\sffamily\color{ForestGreen!70!black},
mdframed={style=mdgreenbox},headpunct={ --- },]{thmgreenbox}

\mdfdefinestyle{mdorangebox}{linecolor=RedOrange,backgroundcolor=RedOrange!5,
	linewidth=2pt,rightline=false,leftline=true,topline=false,bottomline=false,}
\declaretheoremstyle[headfont=\bfseries\sffamily\color{RedOrange!70!black},
mdframed={style=mdorangebox},headpunct={ --- },]{thmorangebox}

\mdfdefinestyle{mdblackbox}{linecolor=black,backgroundcolor=RedViolet!5!gray!5,
	linewidth=3pt,rightline=false,leftline=true,topline=false,bottomline=false,}
\declaretheoremstyle[mdframed={style=mdblackbox}]{thmblackbox}

\declaretheoremstyle[spaceabove=6pt,spacebelow=6pt]{basehead}

\declaretheorem[style=basehead,name=Problem]{problem}
\declaretheorem[style=basehead,name=Problem,numbered=no]{problem*}
\declaretheorem[style=thmbloobox,name=Setup]{setup} % for config geo
\declaretheorem[style=thmredbox,name=Example,sibling=problem]{example}
\declaretheorem[style=thmredbox,name=$\bigstar$ Problem,sibling=problem]{niceprob}
\declaretheorem[style=thmbluebox,name=Theorem,sibling=problem]{theorem}
\declaretheorem[style=thmbluebox,name=Corollary,sibling=theorem]{corollary}
\declaretheorem[style=thmbluebox,name=Proposition,sibling=theorem]{prop}
\declaretheorem[style=thmbluebox,name=Lemma,numberwithin=problem]{lemma}
\declaretheorem[style=thmbluebox,name=Theorem,numbered=no]{theorem*}
\declaretheorem[style=thmbluebox,name=Lemma,numbered=no]{lemma*}
\declaretheorem[style=thmgreenbox,name=Claim,numberwithin=problem]{claim}
\declaretheorem[style=thmgreenbox,name=Claim,numbered=no]{claim*}
\declaretheorem[style=thmblackbox,name=Remark,numberwithin=problem]{remark}
\declaretheorem[style=thmblackbox,name=Remark,numbered=no]{remark*}
\declaretheorem[style=thmgreenbox,name=Definition,sibling=theorem]{definition}
\declaretheorem[style=thmgreenbox,name=Definition,numbered=no]{definition*}
\declaretheorem[style=thmorangebox,name=Warning,sibling=theorem]{warning}
\declaretheorem[style=thmorangebox,name=Warning,numbered=no]{warning*}
\declaretheorem[style=thmorangebox,name=Note,numbered=no]{note}
\declaretheorem[style=thmblackbox,name=Exercise,sibling=problem]{exercise}
\declaretheorem[style=thmblackbox,name=Exercise,numbered=no]{exercise*}
\declaretheorem[style=thmblackbox,name=Question,sibling=problem]{question}
\declaretheorem[style=thmblackbox,name=Question,numbered=no]{question*}

\addtolength{\textheight}{3.14cm}
\providecommand{\ol}{\overline}
\providecommand{\ul}{\underline}
\providecommand{\eps}{\varepsilon}
\providecommand{\half}{\frac{1}{2}}
\providecommand{\dang}{\measuredangle} %% Directed angle
\providecommand{\CC}{\mathbb C}
\providecommand{\FF}{\mathbb F}
\providecommand{\NN}{\mathbb N}
\providecommand{\QQ}{\mathbb Q}
\providecommand{\RR}{\mathbb R}
\providecommand{\ZZ}{\mathbb Z}
\providecommand{\dg}{^\circ}
\providecommand{\ii}{\item}
\providecommand{\setdiff}{\smallsetminus}
\providecommand{\alert}[1]{{\sffamily\textbf{\textcolor{blue}{#1}}}}
\providecommand{\inv}{^{-1}}
\providecommand{\mb}[1]{\mathbf{#1}}
\providecommand{\dif}{\;d}

\reversemarginpar

\newenvironment{soln}{\begin{proof}[Solution]}{\end{proof}}
\newenvironment{walkthrough}{\noindent \textbf{Walkthrough.}}{\medskip}
\newlist{walk}{enumerate}{3}
\setlist[walk]{label=\bfseries (\alph*)}

\providecommand{\pow}{\operatorname{Pow}}
\providecommand{\vocab}[1]{{\textbf{\textcolor{ForestGreen}{#1}}}}
\providecommand{\lcm}{\operatorname{lcm}}
\providecommand{\abs}[1]{\left|#1\right|}

\usepackage{asymptote}
\begin{asydef}
	defaultpen(fontsize(10pt));
	unitsize(1cm);
	size(8cm); // set a reasonable default
	usepackage("amsmath");
	usepackage("amssymb");
	settings.tex="pdflatex";
	settings.outformat="pdf";
	// Replacement for olympiad+cse5 which is not standard
	import geometry;
	// recalibrate fill and filldraw for conics
	void filldraw(picture pic = currentpicture, conic g, pen fillpen=defaultpen, pen drawpen=defaultpen)
	{ filldraw(pic, (path) g, fillpen, drawpen); }
	void fill(picture pic = currentpicture, conic g, pen p=defaultpen)
	{ filldraw(pic, (path) g, p); }
	// some geometry
	pair foot(pair P, pair A, pair B) { return foot(triangle(A,B,P).VC); }
	pair orthocenter(pair A, pair B, pair C) { return orthocentercenter(A,B,C); }
	pair centroid(pair A, pair B, pair C) { return (A+B+C)/3; }
	// cse5 abbreviations
	path CP(pair P, pair A) { return circle(P, abs(A-P)); }
	path CR(pair P, real r) { return circle(P, r); }
	pair IP(path p, path q) { return intersectionpoints(p,q)[0]; }
	pair OP(path p, path q) { return intersectionpoints(p,q)[1]; }
	path Line(pair A, pair B, real a=0.6, real b=a) { return (a*(A-B)+A)--(b*(B-A)+B); }
	// cse5 more useful functions
	picture CC() {
		picture p=rotate(0)*currentpicture;
		currentpicture.erase();
		return p;
	}
	pair MP(Label s, pair A, pair B = plain.S, pen p = defaultpen) {
		Label L = s;
		L.s = "$"+s.s+"$";
		label(L, A, B, p);
		return A;
	}
	pair Drawing(Label s = "", pair A, pair B = plain.S, pen p = defaultpen) {
		dot(MP(s, A, B, p), p);
		return A;
	}
	path Drawing(path g, pen p = defaultpen, arrowbar ar = None) {
		draw(g, p, ar);
		return g;
	}
\end{asydef}

\usepackage{mathrsfs}

\ihead{\footnotesize MAT 314 HW06}
\ohead{\footnotesize Last compiled \today}

\title{\vspace{-2em}MAT 314 HW06}
\author{\large Aditya Pahuja}
\date{\large Due 04/17}
\begin{document}
\maketitle

\begin{problem*}1.2.
	Let $V$ be an abelian group. Prove that if $V$ has a structure of
	$\QQ$-module with its given law of composition as addition,
	then that structure is uniquely determined.
\end{problem*}

\begin{soln}
	We wish to show that there is only one map $\QQ\times V\to V$
	given by $(q,x)\mapsto qx$ such that $1x = x$, $q(x+y) = qx+qy$,
	$(q+r)x = qx + rx$, and $q(rx) = qr(x)$ for $q,r\in\QQ$ and $x,y\in V$.
	
	We have established that there is only one such map $\ZZ\times V\to V$
	(i.e. abelian groups have exactly the natural $\ZZ$-module structure),
	so we want to show that this structure can only be extended to
	that of a $\QQ$-module in at most one way.
	We can see that for $p/q\in\QQ$ with $p,q\in\ZZ$ and $q>0$,
	\[\underbrace{\frac pqx + \cdots + \frac pqx}_{q\text{ times}}
	= q\cdot\frac pqx = px,\]
	so $\frac1q x$ is an element of $V$ that, when added to itself $q$ times,
	equals $x$. If there are two such elements $y$ and $y'$,
	then $q(y-y') = x-x = 0$, so multiplying by $\frac1q$ shows that $y=y'$.
	Thus there can only be at most one choice of $\frac1qx$,
	and then $\frac pqx$ is uniquely determined: it is
	either $0$ (if $p=0$), $\frac 1qx$ added to itself $p$ times
	(if $p>0$), or $-\frac1qx$ added to itself $p$ times (if $p<0$).
\end{soln}

\begin{problem*}2.1.
	Let $R = \CC[x,y]$, and let $M$ be the ideal of $R$
	generated by the two elements $x$ and $y$.
	Is $M$ a free $R$-module?
\end{problem*}

\begin{soln}
	Let $\{P_1, \dots, P_n\}$ be a generating set for $M$.
	We note that $n$ cannot equal $1$ because the ideal generated by
	$P\in R$ has every element divisible by $P$, which implies
	$P\mid x$ and $P\mid y$ so $P$ is a constant polynomial,
	which is impossible because clearly $M$ does not contain
	any polynomials with nonzero constant terms.
	
	 Now we show that $P_1$, \dots, $P_n$ are not independent: indeed,
	 \[-P_2\cdot P_1 + P_1\cdot P_2
	 + 0\cdot P_3 + \cdots + 0\cdot P_n = 0\]
	 but the coefficients of the $P_i$ are not all zero.
	 Therefore no generating set of $M$ is a basis for $M$,
	 so $M$ cannot be a free $R$-module.
\end{soln}

\begin{problem*}3.3.
	Let $A$ and $B$ be $m\times m$ and $n\times n$ matrices, respectively.
	Prove that the trace of the linear operator $f(M) = AMB$
	on the space $R^{m\times n}$ is the product
	$(\operatorname{trace}A)(\operatorname{trace}B)$.
\end{problem*}

\begin{soln}
	Let $E_{i,j}$ be the $m\times n$ matrix whose only nonzero entry
	is a $1$ in row $i$, column $j$.
	Clearly the $E_{i,j}$ are a basis for $R^{m\times n}$.
	The trace of $f$ can be computed as
	\[\sum_{i=1}^m\sum_{j=1}^n \langle E_{i,j}, f(E_{i,j})\rangle\]
	where $\langle X, Y\rangle = \sum_{i,j} X_{i,j}Y_{i,j}$
	for $X,Y\in R^{m\times n}$. The summand is therefore
	the row $i$, column $j$ entry of $AMB$. Now we can see easily that
	$MB$ is the $m\times n$ matrix whose rows are all zeroes
	except for its $i$th row being equal to the $j$th row of $B$,
	and then the entry of $A(MB)$ in row $i$, column $j$ is
	the product $a_{i,i}\cdot b_{j,j}$,
	since column $j$ of $MB$ has $b_{j,j}$ as its only nonzero entry,
	appearing in row $i$, and then the $i$th element in row $i$
	of $A$ is $a_{i,j}$. Therefore the trace is equal to
	\[\sum_{i=1}^m\sum_{j=1}^n a_{i,i}b_{j,j} =
	\left(\sum_{i=1}^ma_{i,i}\right)\left(\sum_{j=1}^na_{j,j}\right).\]
	(The question asked to use the permanence of identities thing,
	but this last equation just directly holds over arbitrary commutative rings
	so I don't really need to appeal to $\CC$ here?)
\end{soln}

\begin{problem*}4.6.
	Let $\varphi\colon\ZZ^k\to\ZZ^k$ be a homomorphism
	given by multiplication by an integer matrix $A$.
	Show that the image of $\varphi$ is of finite index
	if and only if $A$ is nonsingular and that
	if so, then the index is equal to $|\det A|$.
\end{problem*}
 
\begin{soln}
	We wish to show that $\ZZ^k/A\ZZ^k$ is finite with order $|\det A|$
	if and only if $A$ is nonsingular. We can write $A = Q\inv A'P$
	for invertible $P,Q$ where $A'$ is diagonal,
	so that $\ZZ^k/A\ZZ^k\cong\ZZ^k/A'\ZZ^k$. Let the nonzero entries
	of the main diagonal of $A'$ be positive integers $d_1$, \dots, $d_r$
	for $r\le k$ (i.e. Smith normal form, but we don't care about
	the divisibility information here). We can see that
	$A'\ZZ^k = d_1\ZZ\oplus\cdots\oplus d_r\ZZ$,
	so if $A$ (hence $A'$) is nonsingular, then $r = k$, meaning
	\[\ZZ^k/A'\ZZ^k \cong \bigoplus_{i=1}^k \ZZ/d_i\ZZ,\]
	and then the order of this group is $d_1\cdots d_k = \det A' = |\det A|$.
	Conversely if $A$ (hence $A'$) is singular, then $r<k$, so
	\[\ZZ^k/A'\ZZ^k\cong\ZZ^{k-r}\oplus\ZZ/d_1\ZZ\oplus\cdots\oplus \ZZ/d_r\ZZ\]
	which is infinite.
\end{soln}

\begin{problem*}7.4.
	In each case, identify the abelian group
	that has the given presentation matrix:
	\[\begin{pmatrix}
		2\\1
	\end{pmatrix}, \begin{pmatrix}
		0\\5
	\end{pmatrix},
\begin{pmatrix}
	2&0&0
\end{pmatrix},
\begin{pmatrix}
	1&0\\0&1\\0&0
\end{pmatrix},
\begin{pmatrix}
	2&3\\1&2
\end{pmatrix},
\begin{pmatrix}
	2&4\\1&4
\end{pmatrix},
\begin{pmatrix}
	2&4\\6&4
\end{pmatrix},\\
\begin{pmatrix}
	4&6\\2&3
\end{pmatrix}.\]
\end{problem*}

\begin{soln}
	In general for an $m\times n$ matrix $A$,
	the abelian group with presentation matrix $A$ is $\ZZ^m/A\ZZ^n$.
	\begin{itemize}
		\ii The first group is
		\[\ZZ^2/\begin{pmatrix}
			2\\1
		\end{pmatrix}\ZZ = \ZZ^2/\{(2x,x) : x\in\ZZ\} \cong \ZZ,\]
		noting that any element of $\ZZ^2$ can be written as
		$(a + 2x, x)$ for integers $a,x$.
		\ii The second group is
		$\ZZ^2/\{(0, 5x) : x\in\ZZ\}\cong\ZZ\oplus\ZZ/5\ZZ$.
		\ii The third group is $\ZZ/(2\;0\;0)\cdot\ZZ^3 = \ZZ/2\ZZ$.
		\ii The fourth group is the trivial group because
		$A\ZZ^3 = \ZZ^2$, as $A(x,y,0) = (x,y)$ for any $(x,y)\in\ZZ^2$.
		\ii The fifth group is also the trivial group because
		$\det A = 1$, hence it is an automorphism of $\ZZ^2$,
		i.e. $A\ZZ^2 = \ZZ^2$.
		\ii After applying a change of generators and relations,
		we diagonalize this matrix as
		$\begin{pmatrix}
			1&0\\0&4
		\end{pmatrix}$,
		so the corresponding group is
		$\ZZ^2/(\ZZ\oplus 4\ZZ) = \ZZ/4\ZZ$.
		\ii Similar to the previous group, the Smith normal form here is
		$\begin{pmatrix}
			2&0\\0&8
		\end{pmatrix}$, so the quotient is $\ZZ/2\ZZ\oplus\ZZ/8\ZZ$.
		\ii Finally, for the eighth group we can compute that
		$A\ZZ^2 = \{(2x,x) : x\in\ZZ\}$, so the quotient $\ZZ^2/A\ZZ^2$
		is the same as the first group, namely $\ZZ$.
	\end{itemize}
\end{soln}












\end{document}